\documentclass[11pt,letterpaper]{article}

%---------- Packages ----------
\usepackage[margin=0.8in]{geometry}
\usepackage{enumitem}
\usepackage[hidelinks]{hyperref}
\usepackage{titlesec}
\usepackage{array}
\usepackage{setspace}
\usepackage{parskip}

\setlength{\parskip}{3pt}
\setlength{\parindent}{0pt}

%---------- Section formatting ----------
\titleformat{\section}{\bfseries\large}{}{0em}{}[\titlerule]
\titlespacing*{\section}{0pt}{6pt}{3pt}

\titleformat{\subsection}{\bfseries}{}{0em}{}
\titlespacing*{\subsection}{0pt}{4pt}{2pt}

\renewcommand{\labelitemi}{$\bullet$}
\setlist[itemize]{leftmargin=*,itemsep=2pt,topsep=2pt}

%---------- Name ----------
\begin{document}

\begin{center}
    {\LARGE \textbf{Piter Z. Garcia Bautista}}\\[4pt]
    NSF Noyce Scholar \,|\, Data Science \& Bioinformatics Graduate Student\\[3pt]
    Rochester, NY Area \,|\, +1 (646) 288-3300 \,|\, 
    \href{mailto:pzg8794@rit.edu}{pzg8794@rit.edu} \,|\, 
    \href{mailto:piter.garcia@rochester.edu}{piter.garcia@rochester.edu}\\[2pt]
    \href{https://linkedin.com/in/garciapiterz}{linkedin.com/in/garciapiterz} \,|\, 
    \href{https://garciapiterz.com}{garciapiterz.com}
\end{center}

%========== RESEARCH SUMMARY ==========
\section{Research Summary}

NSF Robert Noyce Scholar and PhD-track graduate student specializing in equitable bioinformatics and fairness-aware machine learning for healthcare and educational applications. Developer of the interdisciplinary \textbf{EQUITAS} framework, integrating computational rigor, educational equity, and lived experience to address diagnostic disparities with a focus in culturally and linguistically diverse populations.

%========== EDUCATION ==========
\section{Education}

\textbf{Rochester Institute of Technology (RIT)}, Golisano College of Computing, Rochester, NY\\
\textit{Master of Science in Data Science}, Concentration: Bioinformatics \& Artificial Intelligence \hfill Expected Aug 2026\\
Current GPA: 3.9+ (Excellent Academic Standing)\\
Relevant Coursework: DSCI-633 Foundations of Data Science; DSCI-644 Software Engineering for Data Science; DSCI-601 Applied Data Science I; BIOL-630 Bioinformatics Algorithms; ISTE-780 Data-Driven Knowledge Discovery.

\vspace{4pt}
\textbf{University of Rochester}, Warner School of Education, Rochester, NY\\
\textit{M.S. in Teaching Computer Science K–12} \hfill Expected Aug 2026\\
NSF Robert Noyce Teacher Scholarship Program (Full Funding: \$30,000).\\
Advanced Certificate in Leadership in Disability and Inclusive Practices (in progress).\\
Relevant Coursework: ED-447 Disability and Schools; ED-405 Assessment in Instructional Contexts; ED-452 Instructional Strategies for Inclusive Classrooms; EDE-498 Integrating CS Standards.\\
Service Commitment: 2 years teaching in high-need school districts.

\vspace{4pt}
\textbf{Rochester Institute of Technology (RIT)}, Golisano College of Computing, Rochester, NY\\
\textit{Master of Science in Computer Science} \hfill 2015\\
Concentrations: Computer Graphics, Artificial Intelligence, Big Data Analytics.\\
Research: Technology for Automatic Sign Language Interpretation; Big Data Medical Diagnosis.

\vspace{4pt}
\textbf{Farmingdale State College}, Farmingdale, NY\\
\textit{Bachelor of Science in Computer Engineering} \hfill 2012\\
\textit{Dual Degree in Electrical Engineering Technology \& Computer Engineering Technology}.\\
Capstone: Smart Meter Home Application. Honors: CSTEP Merit Award, Dual BS Scholarship.

%========== CURRENT RESEARCH EXPERIENCE ==========
\section{Research Experience}

\textbf{Lead Researcher – EQUITAS Framework for Diagnostic Justice} \hfill 2024–Present\\
Independent Research Initiative, RIT \& University of Rochester, Rochester, NY
\begin{itemize}
    \item Developed an interdisciplinary framework integrating RNA bioinformatics, fairness-aware machine learning, and trauma-informed care to evaluate and mitigate diagnostic disparities.
    \item Designed equity-focused evaluation pipelines that revealed significant algorithmic bias in existing diagnostic models and guided debiasing strategies for underrepresented groups.
    \item Led cross-institutional collaboration spanning computer science, education, and health sciences, coordinating research tasks and disseminating results as reports and capstone projects.
\end{itemize}

\textbf{Principal Investigator – Equitable Bioinformatics Project (ISTE-780)} \hfill 2024–Present\\
Data Science Program, RIT, Rochester, NY
\begin{itemize}
    \item Pioneered the first systematic application of fairness assessment to RNA secondary structure prediction, integrating AI fairness metrics into bioinformatics workflows.
    \item Implemented a modular bias detection framework identifying critical equity violations (e.g., disparate error rates across demographic groups) and enabling reproducible evaluation.
    \item Developed a 900+ line, well-documented codebase for equity-aware algorithm comparison, achieving statistically significant bias detection across demographic subgroups (\(p < 0.01\)).
\end{itemize}

\textbf{Research Developer – RNA Structure Prediction \& Visualization (BIO-630)} \hfill 2024\\
Bioinformatics Program, RIT, Rochester, NY
\begin{itemize}
    \item Extended single-sequence RNA secondary structure prediction by integrating the Nussinov algorithm with minimum free energy (MFE) enhancements.
    \item Combined dynamic programming with Turner thermodynamic parameters and ViennaRNA tools to improve structure accuracy and feature extraction.
    \item Achieved high accuracy on synthetic and biological motifs and released a validated codebase for high-throughput feature analysis.
\end{itemize}

%========== TEACHING & MENTORING ==========
\section{Teaching \& Mentoring Experience}

\textbf{STEM Tutor \& Academic Mentor} \hfill 2021–Present\\
Varsity Tutors \& Independent Practice (Remote \& In-Person)
\begin{itemize}
    \item Provide comprehensive tutoring in Computer Science, Data Science, Mathematics, and Statistics at undergraduate and graduate levels.
    \item Develop culturally responsive, accessibility-focused lesson plans supporting diverse learners, including students with disabilities and neurodivergent students.
    \item Mentor underrepresented students in STEM, maintaining a high satisfaction rate and documented improvement in academic performance and confidence.
\end{itemize}

\textbf{Faculty Assistant – Graduate Office} \hfill 2013–2015\\
Golisano College of Computing, RIT, Rochester, NY
\begin{itemize}
    \item Supported program coordination, student services, and graduate operations, serving as a frontline resource for graduate student advising and documentation.
    \item Assisted faculty and staff with communications, event logistics, and administrative workflows that improved student experience and program efficiency.
\end{itemize}

\textbf{Private Tutor – Computer Science} \hfill 2013–2015\\
Independent Practice, Rochester, NY
\begin{itemize}
    \item Delivered targeted tutoring in Java, algorithms, and theoretical CS foundations for undergraduate and graduate students.
    \item Designed individualized learning plans to support students preparing for exams, projects, and research-intensive coursework.
\end{itemize}

%========== PROFESSIONAL EXPERIENCE ==========
\section{Industry Experience}

\textbf{Data Solutions Engineer} \hfill Sep 2022–Aug 2024\\
VEDADATA Inc., New York, NY
\begin{itemize}
    \item Designed and implemented data-driven solutions using Python, Pandas, and cloud infrastructure to address complex business and analytics problems.
    \item Architected scalable, maintainable pipelines with an emphasis on accessibility, robustness, and reproducibility.
    \item Collaborated with cross-functional teams to deliver production-grade software aligned with stakeholder requirements.
\end{itemize}

\textbf{AI Data Solutions Engineer} \hfill Jan 2021–Sep 2022\\
VIOME Inc., New York, NY
\begin{itemize}
    \item Developed AI-driven healthcare solutions for personalized health recommendations, focusing on equitable outcomes and responsible data use.
    \item Implemented best practices for AWS deployment, data management, and SDLC compliance in regulated environments.
    \item Led development of predictive models integrating multi-omics and clinical data to support diagnostic decision-making.
\end{itemize}

\textbf{Computer \& Data Science Consultant} \hfill Mar 2017–Jan 2021\\
Samasta Health Foundation (Remote/Hybrid)
\begin{itemize}
    \item Provided pro-bono consulting to architect smart systems supporting children in need in low-resource settings.
    \item Developed AI-driven healthcare tools emphasizing equity, accessibility, and trauma-informed design.
    \item Collaborated with healthcare professionals to design solutions aligned with community needs and diagnostic equity goals.
\end{itemize}

%========== PUBLICATIONS & ACADEMIC OUTPUT ==========
\section{Selected Publications \& Academic Contributions}

\begin{itemize}
    \item Garcia, P. (2025). \textit{Equitable Bioinformatics: Enhancing Diagnostic Decision-Making through RNA and Biomarker Data}. ISTE-780 Project Report, Rochester Institute of Technology.
    \item Garcia, P. (2025). \textit{Trauma-Informed Bioinformatics for Diagnostic Justice: An EQUITAS Framework for Integrative Healthcare and Educational Systems}. ED-447 Capstone Project, University of Rochester.
    \item Garcia, P. (2025). \textit{RNA Secondary Structure Prediction and Visualization: Implementing the Nussinov Algorithm with MFE Enhancements}. BIO-630 Final Project, Rochester Institute of Technology.
    \item Garcia, P. (2015). \textit{ASL Recognition Application: Image Processing for Sign Language Translation}. M.S. Thesis, Rochester Institute of Technology.
    \item Garcia, P. (2014). \textit{Big Data Medical Diagnosis: A Three-Method Approach to Improve Diagnostic Accuracy}. RIT Graduate Research Symposium Presentation.
\end{itemize}

%========== TECHNICAL SKILLS ==========
\section{Technical Skills}

\textbf{Programming Languages:} Python, R, Java, C++, C\#, MATLAB, SQL, JavaScript, HTML/CSS, Bash.\\
\textbf{Bioinformatics \& Computational Biology:} ViennaRNA, Nussinov Algorithm, Turner Parameters, BLAST, SHAP Analysis.\\
\textbf{Machine Learning \& AI:} TensorFlow, Keras, scikit-learn, XGBoost, PyTorch, LSTM, fairness-aware ML.\\
\textbf{Data Science \& Analytics:} Pandas, NumPy, SciPy, Matplotlib, Plotly, Seaborn, SPSS, statistical analysis.\\
\textbf{Cloud \& Infrastructure:} AWS, GCP, Azure, Docker, Git/GitHub, Linux/Unix systems.\\
\textbf{Databases:} PostgreSQL, MySQL, MongoDB, NoSQL.\\
\textbf{Specialized Frameworks:} EQUITAS fairness assessment, bias detection, disparate impact analysis.

%========== AWARDS & HONORS ==========
\section{Awards \& Honors}

\begin{itemize}
    \item NSF Robert Noyce Teacher Scholarship (2024–2026) – Full funding for M.S. in Teaching Computer Science K–12.
    \item MS International Scholarship (2012–2015), MESCYT \& RIT – Academic excellence recognition.
    \item Honorable Mention Technology Award (2012), NYS STEP/CSTEP – 2nd place capstone project.
    \item IEEE Alumni Member (2009–Present) – Member No. 90613000.
    \item Academic Excellence Merit (2002–2006), Mateo Lara High School.
\end{itemize}

%========== FACULTY COLLABORATION INTERESTS ==========
\section{Faculty Collaboration Interests}

Interested in collaborating with faculty whose work aligns with equitable AI, bioinformatics, and computer science education, including:
\begin{itemize}
    \item Dr. Qi Yu – Human-centered AI \& machine learning; inspiration for equitable bioinformatics project development.
    \item Dr. Fernando Rodriguez – Comparative genomics \& molecular evolution; foundational guidance for RNA structure work.
    \item Dr. Travis Desell – Neuroevolution \& distributed computing; Data Science MS program advisor and algorithmic mentor.
    \item Dr. Feng Cui – Computational biology \& AI in medicine; alignment with health equity-focused bioinformatics.
    \item Dr. Evan Selinger – AI ethics \& algorithmic accountability; ethical foundation for the EQUITAS framework.
    \item Dr. Christian Newman – Software engineering \& educational technology; supports technical and teaching development.
\end{itemize}

\end{document}
